% Ad Atticum 15.5 --- headnote
% ad-atticum-15-05, 28 May 44 BC, in Tusculano

\textit{Cicero to Atticus, written at the Tusculan villa on 28
May 44 BC --- Perseus dateline \textit{Scr. in Tusculano v K.
Iun. a. 710 (44)}. The Liberators' courier has just returned
with letters from both Brutus and Cassius. Brutus is pressing
Cicero for counsel between two courses (the province assigned
him by the Senate's decree, or some bolder step); Cassius is
pressing him to make Hirtius --- the consul-designate and
Cicero's own former pupil --- as well-disposed to their cause
as possible. Cicero registers blank dismay: he has nothing to
write and means to keep silence. Cassius's request in
particular he dismisses with the daggered Greek tag
\textit{[Greek: ho th\=esauros anthrakes]} (``the treasure is
coals''), proverbial for hopes that turn out worthless ---
Hirtius, however he may have been earlier, will not be made
better by Cicero's authority now.}

\textit{Section 2 reports that Balbus and Oppius confirm the
provincial decree Atticus had mentioned; Hirtius is already
at his own Tusculan villa and urges Cicero to keep away from
Rome, on grounds of danger. Cicero's own reason is different:
not Antony's suspicion (he is, if anything, willing to be seen
displeased at Antony's good fortune), but simply that he does
not want to have to look at the man. Section 3 turns to the
disbanded veterans, who according to a letter to Varro from an
unnamed source are talking in the most disgraceful way, so
that any dissenter in Rome will be in real peril; and to
Lucius Antonius moving against Decimus Brutus. The famous
formulation closes the letter --- a city in which Cicero once
``not only flourished at the height, but also served in
slavery with some measure of dignity'' --- and he resolves at
least, for now, not to return to Rome, leaving the question of
whether to leave Italy altogether for Atticus to decide with
him face to face.}
